\chapter{Considerações Finais}

\section{Conclusão}

Ao longo do trabalho foram aplicados conceitos fundamentais da engenharia de \textit{software}, desde o levantamento da literatura e a definição de requisitos até a modelagem, a arquitetura e a prototipação de alta fidelidade.
Conforme dito anteriormente, o objetivo do projeto é o desenvolvimento de um aplicativo \textit{mobile}, voltado à organização da rotina de estudos de estudantes, por meio da integração do cronômetro Pomodoro configurável, \textit{flashcards} com repetição espaçada, gamificação e \textit{check-ins} de humor em uma única plataforma.

O levantamento da literatura realizado, com consultas na base científica Scopus, fundamenta a seleção das funcionalidades propostas com base em evidências acadêmicas sobre gestão do tempo, repetição espaçada, gamificação e aprendizagem autorregulada; essa fundamentação teórica, no entanto, embasa a escolha das funcionalidades, e não substitui a validação de sua implementação, usabilidade e efetividade, que deverá ocorrer durante e após a etapa de desenvolvimento.
A literatura consultada indica que a gestão do tempo e a autoavaliação estão entre os fatores mais associados ao desempenho acadêmico, o que motivou a priorização do cronômetro Pomodoro e do algoritmo de repetição espaçada.
Da mesma forma, os estudos sobre gamificação e aprendizagem autorregulada sustentam, respectivamente, os elementos de engajamento e os \textit{check-ins} de bem-estar emocional como recursos relevantes, ainda que a gamificação, conforme discutido na Seção~\ref{sec:gamificacao}, tenda a impactar sobretudo a motivação do estudante.

A arquitetura em camadas planejada, com \textit{backend} em Node.js e Express, \textit{frontend} em React Native com Expo e banco de dados relacional PostgreSQL, foi projetada com o objetivo de favorecer a escalabilidade, a manutenibilidade e a conformidade com a LGPD, propriedades que também dependerão das decisões tomadas durante a implementação e a fase de testes.
A condução do \textit{framework} \textit{Lean Inception}, com a participação de duas proto-personas e do autor, contribuiu para explorar e organizar hipóteses sobre as necessidades dos usuários, resultando em um \textit{Canvas} MVP e em um sequenciador de funcionalidades que orientam o desenvolvimento incremental; por se tratar de uma amostra reduzida de participantes, essas hipóteses ainda carecem de validação junto a uma base maior de usuários reais.
A aplicação combinada das metodologias \textit{Scrum}, Kanban, XP e TDD busca apoiar um ciclo de desenvolvimento ágil, rastreável e orientado à qualidade.

Portanto, esta etapa do projeto entrega uma proposta teoricamente fundamentada e centrada no usuário, incluindo a definição de um escopo de MVP reduzido e rastreável até os requisitos funcionais, servindo de base para a implementação a ser conduzida na próxima etapa.

\section{Trabalhos Futuros}
\label{sec:trabalhos_futuros}

Diante dos resultados alcançados, torna-se viável a continuidade do desenvolvimento do projeto, no qual serão implementadas as funcionalidades priorizadas no Sequenciador e no \textit{Canvas} MVP, seguindo as ondas de desenvolvimento definidas durante o \textit{Lean Inception}.
Por se tratar de um projeto gratuito e \textit{open source}, distribuído sob a licença MIT no repositório \url{https://github.com/H0lzz/TCC1.2}, incrementos futuros podem ser propostos e implementados por outros desenvolvedores interessados em contribuir com o aplicativo.

Como trabalhos futuros, destacam-se:

\begin{itemize}
    \item a validação das proto-personas e das hipóteses de necessidade levantadas neste trabalho junto a uma amostra maior de usuários reais, por meio de entrevistas, testes de usabilidade ou pesquisas de campo;
    \item a implementação de um modo \textit{offline-first}, com armazenamento local e sincronização posterior dos dados, funcionalidade identificada como desejável durante o \textit{Lean Inception}, mas classificada fora do escopo do MVP em razão de seu esforço técnico;
    \item a incorporação de Inteligência Artificial para a geração automatizada de \textit{flashcards} a partir de materiais enviados pelo usuário, como PDFs e imagens de anotações, funcionalidade que não faz parte do escopo do MVP, mas que poderia representar um diferencial relevante;
    \item a implementação de todas as funcionalidades classificadas como incrementos do sequenciador.
\end{itemize}
